% ============================================================
% UNIDAD 1: DERECHO AERONÁUTICO Y ROL REGULADOR DEL ESTADO
% ============================================================
\clearpage
\section*{UNIDAD 1: DERECHO AERONÁUTICO}
\subsection*{Marco internacional, marco nacional y rol regulador del Estado}

\begin{recuadro}
\textbf{Objetivo de la unidad:} conocer de dónde salen las reglas que rigen la aviación, quién
las dicta y quién controla que se cumplan, tanto a nivel internacional (OACI) como nacional
(Código Aeronáutico, ANAC, JST, RAAC), y ubicar el lugar del técnico aeronáutico dentro de ese
sistema.
\end{recuadro}

% ------------------------------------------------------------
\section*{1. ¿QUÉ ES EL DERECHO AERONÁUTICO?}

El \textbf{derecho aeronáutico} es el conjunto de normas, nacionales e internacionales, que
regulan la actividad aeronáutica: la circulación de aeronaves, la infraestructura (aeródromos,
servicios de navegación aérea), las aeronaves y su aeronavegabilidad, el personal aeronáutico y
sus licencias, el transporte de pasajeros y carga, la responsabilidad por daños y la
investigación de accidentes.

\subsubsection*{1.1. Características}
\begin{itemize}
  \item \textbf{Internacional:} un avión cruza fronteras en minutos; las reglas tienen que ser
        parecidas en todos los países. Por eso gran parte de las normas nacionales copian las
        normas de la OACI.
  \item \textbf{Dinámico:} la tecnología cambia rápido (motores a reacción, GPS, drones) y las
        normas se actualizan constantemente.
  \item \textbf{Autónomo:} tiene principios propios, distintos del derecho marítimo o del
        terrestre.
  \item \textbf{Integral:} mezcla derecho público (Estado, autoridad, sanciones) y privado
        (contratos de transporte, seguros, responsabilidad).
  \item \textbf{Técnico:} muchas de sus normas son reglas técnicas (cómo se mantiene un avión,
        qué instrumentos debe llevar, cuántas horas puede trabajar un tripulante).
\end{itemize}

\subsubsection*{1.2. Fuentes}
\begin{center}
\begin{tikzpicture}[scale=0.85, every node/.style={font=\small, align=center}]
  \fill[ppAzul!85] (0,4.2) -- (-1.3,3.15) -- (1.3,3.15) -- cycle;
  \node[white] at (0,3.5) {\textbf{1}};
  \fill[ppAzul!65] (-1.3,3.1) -- (1.3,3.1) -- (2.4,2.1) -- (-2.4,2.1) -- cycle;
  \fill[ppAzul!45] (-2.4,2.05) -- (2.4,2.05) -- (3.5,1.05) -- (-3.5,1.05) -- cycle;
  \fill[ppAzul!25] (-3.5,1.0) -- (3.5,1.0) -- (4.6,0) -- (-4.6,0) -- cycle;
  \node at (0,2.6) {\textbf{2}};
  \node at (0,1.55) {\textbf{3}};
  \node at (0,0.5) {\textbf{4}};
  \node[anchor=west] at (1.5,3.6) {Constitución Nacional};
  \node[anchor=west] at (2.6,2.6) {Tratados internacionales (Convenio de Chicago)};
  \node[anchor=west] at (3.7,1.55) {Leyes y decretos (Código Aeronáutico, Ley 27.514)};
  \node[anchor=west] at (4.8,0.5) {Reglamentos (RAAC, resoluciones de la ANAC)};
\end{tikzpicture}
\end{center}
A estas fuentes escritas se suman la \textbf{costumbre} (prácticas aceptadas de la industria),
la \textbf{jurisprudencia} (fallos judiciales) y la \textbf{doctrina} (estudios de juristas).
Además, los documentos técnicos de los fabricantes y de la OACI funcionan en la práctica como
``normas'' de cumplimiento obligatorio cuando una regulación los incorpora.

% ------------------------------------------------------------
\section*{2. EL MARCO INTERNACIONAL}

\subsubsection*{2.1. Antecedentes}
\begin{itemize}
  \item \textbf{Convención de París (1919):} primer gran tratado; reconoció que cada Estado es
        soberano sobre el espacio aéreo que cubre su territorio.
  \item \textbf{Convenio de Varsovia (1929):} reguló la responsabilidad del transportista aéreo
        internacional por daños a pasajeros y equipajes. Hoy lo reemplaza, para la mayoría de
        los países, el \textbf{Convenio de Montreal (1999)}.
  \item \textbf{Convenio de Chicago (1944):} es la ``constitución'' de la aviación civil
        internacional.
\end{itemize}

\subsubsection*{2.2. El Convenio sobre Aviación Civil Internacional (Chicago, 1944)}
Firmado el 7 de diciembre de 1944, mientras terminaba la Segunda Guerra Mundial. La Argentina lo
aprobó por la \textbf{Ley 13.891}. Sus principios más importantes:
\begin{itemize}
  \item \textbf{Soberanía:} cada Estado tiene soberanía plena y exclusiva sobre el espacio aéreo
        situado sobre su territorio (artículo 1).
  \item \textbf{Nacionalidad de las aeronaves:} toda aeronave tiene la nacionalidad del Estado
        en que está matriculada, y solo puede estar matriculada en un Estado.
  \item \textbf{Documentos a bordo:} toda aeronave en vuelo internacional debe llevar, entre
        otros, su \textbf{certificado de matrícula}, su \textbf{certificado de
        aeronavegabilidad} y las \textbf{licencias} de la tripulación.
  \item \textbf{Reconocimiento mutuo:} los certificados de aeronavegabilidad y las licencias
        emitidos por un Estado son reconocidos por los demás, siempre que cumplan como mínimo
        las normas de la OACI.
  \item \textbf{Creación de la OACI}, que comenzó a funcionar en 1947.
\end{itemize}

\paragraph*{Las libertades del aire.}
Son los derechos que los Estados se conceden entre sí para operar vuelos comerciales:
\begin{center}
\small
\begin{tabular}{|c|p{0.80\linewidth}|}
\hline
\textbf{Libertad} & \textbf{Derecho} \\ \hline
1.ª & Sobrevolar el territorio de otro Estado sin aterrizar. \\ \hline
2.ª & Aterrizar en otro Estado con fines no comerciales (combustible, mantenimiento). \\ \hline
3.ª & Llevar pasajeros o carga desde el propio país hacia otro país. \\ \hline
4.ª & Traer pasajeros o carga desde otro país hacia el propio. \\ \hline
5.ª & Llevar pasajeros entre dos países extranjeros, en un vuelo que empieza o termina en el
propio país. \\ \hline
6.ª a 9.ª & Combinaciones más amplias, incluido el \textbf{cabotaje} (vuelos internos en otro
país), que muy pocos Estados conceden. \\ \hline
\end{tabular}
\end{center}

\subsubsection*{2.3. La OACI}
La \textbf{Organización de Aviación Civil Internacional} (OACI, en inglés ICAO) es un organismo
especializado de las Naciones Unidas, con sede en Montreal (Canadá), del que forman parte casi
todos los países del mundo.
\begin{itemize}
  \item \textbf{Asamblea:} todos los Estados; se reúne cada tres años.
  \item \textbf{Consejo:} órgano ejecutivo permanente; adopta las normas.
  \item \textbf{Comisión de Aeronavegación:} estudia y propone las normas técnicas.
  \item \textbf{Secretaría} y \textbf{oficinas regionales} (la Argentina depende de la Oficina
        Sudamericana, en Lima).
\end{itemize}

\paragraph*{SARPs: normas y métodos recomendados.}
La OACI dicta \textbf{SARPs} (\textit{Standards and Recommended Practices}):
\begin{itemize}
  \item \textbf{Norma:} su aplicación uniforme es \textbf{necesaria}. Si un Estado no la puede
        cumplir, debe \textbf{notificar la diferencia} a la OACI.
  \item \textbf{Método recomendado:} su aplicación es \textbf{conveniente}; los Estados procuran
        cumplirlo.
\end{itemize}
Los SARPs se publican en los \textbf{19 Anexos} al Convenio, y se complementan con
\textbf{documentos} (manuales) que explican cómo aplicarlos, como el Doc 9859, \textit{Manual de
gestión de la seguridad operacional}.

\begin{center}
\small
\begin{tabular}{|c|p{0.36\linewidth}||c|p{0.36\linewidth}|}
\hline
\textbf{Anexo} & \textbf{Tema} & \textbf{Anexo} & \textbf{Tema} \\ \hline
1 & \textbf{Licencias al personal} & 11 & Servicios de tránsito aéreo \\ \hline
2 & Reglamento del aire & 12 & Búsqueda y salvamento \\ \hline
3 & Meteorología & 13 & \textbf{Investigación de accidentes e incidentes} \\ \hline
4 & Cartas aeronáuticas & 14 & Aeródromos \\ \hline
5 & Unidades de medida & 15 & Servicios de información aeronáutica \\ \hline
6 & \textbf{Operación de aeronaves} & 16 & Protección del medio ambiente \\ \hline
7 & Marcas de nacionalidad y matrícula & 17 & Seguridad (protección contra actos de
interferencia ilícita) \\ \hline
8 & \textbf{Aeronavegabilidad} & 18 & Mercancías peligrosas \\ \hline
9 & Facilitación & 19 & \textbf{Gestión de la seguridad operacional} \\ \hline
10 & Telecomunicaciones aeronáuticas (incluye radioayudas y comunicaciones: el corazón de la
aviónica) & & \\ \hline
\end{tabular}
\end{center}

\paragraph*{Auditorías.}
La OACI audita a los Estados mediante el \textbf{USOAP} (Programa universal de auditoría de la
vigilancia de la seguridad operacional), que mide si cada país tiene la capacidad de controlar
eficazmente su aviación civil.

\subsubsection*{2.4. Otras autoridades de referencia}
\begin{itemize}
  \item \textbf{FAA} (Administración Federal de Aviación, EE.\,UU.): sus regulaciones (\textbf{FAR},
        14 CFR) son la base histórica de las RAAC, y la mayor parte de la flota argentina es de
        diseño estadounidense.
  \item \textbf{EASA} (Agencia de la Unión Europea para la Seguridad Aérea): autoridad europea;
        certifica a Airbus y a muchos fabricantes de equipos.
  \item \textbf{SRVSOP} (Sistema Regional de Cooperación para la Vigilancia de la Seguridad
        Operacional): sistema de cooperación de los países sudamericanos, que elabora los
        \textbf{LAR} (Reglamentos Aeronáuticos Latinoamericanos) para armonizar las normas de la
        región.
\end{itemize}

\begin{recuadro}
\textbf{Seguridad operacional (\textit{safety}) y seguridad de la aviación (\textit{security}).}
En español usamos ``seguridad'' para dos cosas distintas. La \textbf{seguridad operacional}
(\textit{safety}) trata de evitar accidentes por fallas técnicas, errores humanos o
condiciones peligrosas: es el tema de esta materia (Anexos 6, 8, 13 y 19). La \textbf{seguridad
de la aviación} o protección (\textit{security}, AVSEC) trata de evitar actos de interferencia
ilícita: secuestros, sabotajes, terrorismo (Anexo 17); en la Argentina la ejerce, entre otros,
la Policía de Seguridad Aeroportuaria.
\end{recuadro}

% ------------------------------------------------------------
\section*{3. EL MARCO NACIONAL}

\subsubsection*{3.1. El Código Aeronáutico (Ley 17.285)}
Sancionado en 1967, es la ley que rige la aeronáutica civil en el territorio argentino, su mar
territorial y el espacio aéreo que los cubre. Trata, entre otros temas:
\begin{itemize}
  \item la circulación aérea y la infraestructura (aeródromos, servicios);
  \item las \textbf{aeronaves}: definición, nacionalidad, \textbf{Registro Nacional de Aeronaves},
        matrícula, certificado de aeronavegabilidad;
  \item el \textbf{personal aeronáutico}: necesita una licencia o certificado de idoneidad
        expedido por la autoridad para ejercer sus funciones;
  \item la explotación comercial y el contrato de transporte aéreo;
  \item la responsabilidad por daños a pasajeros, cargas y terceros en la superficie;
  \item la búsqueda, asistencia y salvamento, y la investigación de accidentes;
  \item las faltas, delitos y sanciones.
\end{itemize}

\paragraph*{Reformas recientes.}
Entre 2023 y 2025 el Código fue modificado varias veces por decretos de necesidad y urgencia
(DNU 70/2023, 941/2024 y 338/2025), con el objetivo declarado de desregular la actividad
aerocomercial: apertura del mercado a nuevas empresas y rutas (``cielos abiertos''),
simplificación de trámites, eliminación de requisitos considerados obsoletos y posibilidad de
que la autoridad delegue algunas funciones de certificación en personas habilitadas. También se
calificó a la aeronáutica civil aerocomercial como \textbf{servicio esencial}. Las reformas son
objeto de debate entre quienes destacan la mayor competencia y quienes advierten sobre sus
efectos en el empleo y en la vigilancia estatal; para esta materia lo importante es saber que
\textbf{las normas cambian} y que hay que trabajar siempre con la versión vigente.

\subsubsection*{3.2. Los organismos del sistema aeronáutico argentino}
% Logos: colocar los archivos en logos/organismos/ con estos nombres (png, jpg o pdf):
% transporte, anac, jst, eana, orsna, psa, fuerza_aerea (Fuerza Aérea Argentina), smn.
% Si un archivo no está, se muestra un recuadro con la sigla.
\newcommand{\logoorg}[2]{\vspace{2pt}\logoorgimg{#1}{#2}\vspace{2pt}}
\newcommand{\logoorgimg}[2]{%
  \IfFileExists{logos/organismos/#1.png}{\includegraphics[width=2.9cm,height=1.1cm,keepaspectratio]{logos/organismos/#1.png}}{%
  \IfFileExists{logos/organismos/#1.jpg}{\includegraphics[width=2.9cm,height=1.1cm,keepaspectratio]{logos/organismos/#1.jpg}}{%
  \IfFileExists{logos/organismos/#1.pdf}{\includegraphics[width=2.9cm,height=1.1cm,keepaspectratio]{logos/organismos/#1.pdf}}{%
  \fcolorbox{gray!60}{gray!8}{\parbox[c][1.0cm][c]{2.6cm}{\centering\scriptsize\color{gray}\textbf{#2}\\logo}}}}}}
\newcolumntype{M}[1]{>{\centering\arraybackslash}m{#1\linewidth}}
\newcolumntype{N}[1]{>{\raggedright\arraybackslash}m{#1\linewidth}}
\begin{center}
\footnotesize
\renewcommand{\arraystretch}{1.05}
\begin{tabular}{|M{0.20}|N{0.18}|N{0.52}|}
\hline
\textbf{Logo} & \textbf{Organismo} & \textbf{Función} \\ \hline
\logoorg{transporte}{ST} & \textbf{Secretaría de Transporte} & Define la política aeronáutica nacional (rutas, mercado,
tarifas, acuerdos con otros países). \\ \hline
\logoorg{anac}{ANAC} & \textbf{ANAC} -- Administración Nacional de Aviación Civil & Creada en 2007 (Decreto 239/2007),
es la \textbf{autoridad aeronáutica civil}. Dicta las RAAC, certifica aeronaves, talleres,
operadores y centros de instrucción, otorga las licencias del personal y fiscaliza su
cumplimiento. Su \textbf{Dirección Nacional de Seguridad Operacional (DNSO)} tiene a su cargo
las áreas de aeronavegabilidad, operaciones y licencias. \\ \hline
\logoorg{jst}{JST} & \textbf{JST} -- Junta de Seguridad en el Transporte & Creada por la \textbf{Ley 27.514} (2019),
reemplazó a la JIAAC. Investiga los accidentes e incidentes de transporte (aéreo, ferroviario,
marítimo, automotor) para emitir \textbf{recomendaciones} de seguridad. \textbf{No busca culpables
ni determina responsabilidades}: eso es tarea de la Justicia. \\ \hline
\logoorg{eana}{EANA} & \textbf{EANA} -- Empresa Argentina de Navegación Aérea & Empresa estatal que presta los
servicios de navegación aérea: control de tránsito aéreo, comunicaciones, radioayudas (VOR, ILS,
DME), radares e información aeronáutica. Emplea muchos técnicos en electrónica y aviónica. \\
\hline
\logoorg{orsna}{ORSNA} & \textbf{ORSNA} -- Organismo Regulador del Sistema Nacional de Aeropuertos & Controla las
concesiones de los aeropuertos. \\ \hline
\logoorg{psa}{PSA} & \textbf{PSA} -- Policía de Seguridad Aeroportuaria & Seguridad de la aviación (AVSEC) en los
aeropuertos. \\ \hline
\logoorg{fuerza_aerea}{FAA} & \textbf{Fuerza Aérea Argentina} & Aviación militar y defensa del espacio aéreo. Sus organismos
de mantenimiento, como el Área de Material Quilmes, son parte de la historia de esta escuela. \\
\hline
\logoorg{smn}{SMN} & \textbf{SMN} -- Servicio Meteorológico Nacional & Meteorología aeronáutica. \\ \hline
\end{tabular}
\end{center}

\subsubsection*{3.3. Las RAAC: Regulaciones Argentinas de Aviación Civil}
Son los reglamentos técnicos que dicta la ANAC para cumplir el Código Aeronáutico y los Anexos
de la OACI. Están organizadas en \textbf{Partes} numeradas, con una numeración muy parecida a la
de las FAR de Estados Unidos.

\begin{center}
\small
\begin{tabular}{|c|p{0.76\linewidth}|}
\hline
\textbf{Parte} & \textbf{Contenido} \\ \hline
1 & Definiciones y abreviaturas \\ \hline
11 & Reglas para el desarrollo, aprobación y enmienda de las RAAC \\ \hline
13 & Notificación e investigación de accidentes e incidentes \\ \hline
21 & Certificación de productos aeronáuticos (certificados tipo, de aeronavegabilidad,
aprobación de partes) \\ \hline
23, 25, 27, 29 & Estándares de aeronavegabilidad: aviones normales, de transporte, helicópteros
normales y de transporte \\ \hline
39 & Directivas de aeronavegabilidad \\ \hline
\textbf{43} & \textbf{Mantenimiento, mantenimiento preventivo, reconstrucción y alteraciones} \\
\hline
45 & Identificación de productos y marcas \\ \hline
47 & Registro de aeronaves \\ \hline
61 & Licencias de pilotos \\ \hline
\textbf{65} & \textbf{Personal aeronáutico, excepto tripulantes de vuelo (mecánicos, MERA,
despachantes, controladores)} \\ \hline
67 & Certificación médica aeronáutica \\ \hline
91 & Reglas de vuelo y operación general \\ \hline
100 & Aeronaves pilotadas a distancia (drones, RPAS) \\ \hline
119, \textbf{121}, 135 & Certificación y operación de explotadores aerocomerciales (\textbf{121}:
líneas aéreas regulares; 135: taxis aéreos y operaciones menores) \\ \hline
137 & Trabajo aéreo (aviación agrícola, fotografía, etc.) \\ \hline
141, 142 & Centros de instrucción de pilotos y de entrenamiento \\ \hline
\textbf{145} & \textbf{Talleres Aeronáuticos de Reparación (TAR)} \\ \hline
\textbf{147} & \textbf{Centros de instrucción para mecánicos de mantenimiento} \\ \hline
153, 154, 77 & Aeródromos y obstáculos \\ \hline
\end{tabular}
\end{center}

\paragraph*{Tendencia actual.}
Desde 2024 la ANAC está actualizando muchas Partes (21, 61, 91, 121, 135, 147, 153 y otras) para
simplificarlas y alinearlas con las normas de la FAA y la EASA, por ejemplo reconociendo
directamente certificados tipo emitidos por autoridades extranjeras. Por eso es clave consultar
siempre la última enmienda publicada.

% ------------------------------------------------------------
\clearpage
\section*{4. EL ROL REGULADOR DEL ESTADO}

El Estado tiene la responsabilidad de que la aviación sea segura. Para eso cumple varias
funciones, que conviene no confundir:

\begin{center}
\small
\begin{tabular}{|p{0.22\linewidth}|p{0.42\linewidth}|p{0.24\linewidth}|}
\hline
\textbf{Función} & \textbf{En qué consiste} & \textbf{Quién la ejerce} \\ \hline
Política aeronáutica & Decidir qué mercado, qué rutas, qué acuerdos internacionales. &
Secretaría de Transporte \\ \hline
Normar & Dictar las regulaciones técnicas. & ANAC (RAAC) \\ \hline
Certificar y licenciar & Aprobar aeronaves, talleres, operadores, escuelas y personas. & ANAC \\
\hline
Vigilar (fiscalizar) & Inspeccionar y auditar que se siga cumpliendo después de certificar. &
ANAC \\ \hline
Sancionar & Aplicar multas, suspender o revocar licencias y certificados. & ANAC \\ \hline
Investigar & Determinar las causas de accidentes para prevenir otros. & JST \\ \hline
Prestar servicios & Control de tránsito aéreo, radioayudas, meteorología. & EANA, SMN \\ \hline
\end{tabular}
\end{center}

\paragraph*{Los elementos críticos de un sistema de vigilancia.}
La OACI define ocho \textbf{elementos críticos} que todo Estado debe tener para vigilar
eficazmente su aviación: (1) legislación aeronáutica básica; (2) reglamentos de explotación
específicos; (3) sistema y funciones estatales (una autoridad organizada); (4) personal técnico
calificado e instruido; (5) orientación técnica, instrumentos y suministro de información crítica;
(6) obligaciones de otorgamiento de licencias, certificaciones, autorizaciones y aprobaciones;
(7) obligaciones de vigilancia; (8) solución de problemas de seguridad operacional. Además, cada
Estado debe tener un \textbf{Programa Estatal de Seguridad Operacional (SSP)}, que se estudia en
la Unidad 6.

% ------------------------------------------------------------
\section*{5. ACTIVIDADES DE LA UNIDAD}
\begin{enumerate}
  \item \textbf{Esquema institucional:} dibujar un esquema del sistema aeronáutico argentino con
        la OACI, la Secretaría de Transporte, la ANAC, la JST, la EANA, el ORSNA y la PSA,
        indicando con flechas quién regula, quién certifica, quién investiga y quién presta
        servicios.
  \item \textbf{Cuadro comparativo:} comparar para un mismo tema (por ejemplo, la licencia del
        mecánico, o los talleres de reparación) qué dice la OACI (Anexo), la FAA (FAR), la EASA
        y la ANAC (RAAC).
  \item \textbf{Libertades del aire:} para tres vuelos reales que operan en Ezeiza, indicar qué
        libertad del aire están ejerciendo.
  \item \textbf{Búsqueda normativa:} entrar al sitio de la ANAC, encontrar la RAAC Parte 43 y la
        Parte 145 vigentes y registrar su número de edición y la última enmienda.
  \item \textbf{Debate:} ¿qué ventajas y riesgos tiene armonizar las normas argentinas con las de
        la FAA? ¿Qué pasa con la vigilancia del Estado cuando se simplifican los requisitos?
\end{enumerate}
